%% FullCourse_10Lecs -- Lecture 7, Topic 7.1: Why LLMs + Tokenization + Word Embeddings
%% Source: ./archive_pre_split/Lec07_LLM.tex (sections 1-3)
%% Pass B: layer in refinements from
%%         ../../MiniCourse_8hr/Slides/Module3_LLM_Text/M3_T1_Tokenization_Embeddings.tex

%% Shared preamble for the FullCourse_10Lecs topic decks.
%%
%% Each topic .tex \input's this file, then defines its own \title, \date
%% override (if any), \graphicspath, and \begin{document}.
%%
%% Reconstructed verbatim from the inlined copy preserved in
%% Lec10_Case_Studies/Lec10_T8_Case6_DeepRamsey.tex (lines 23-190), which is
%% explicitly annotated there as "from ../shared_preamble.tex, copied verbatim".
%% Base preamble matches MiniCourse_8hr/Slides/shared_preamble.tex; this version
%% additionally provides \sectiondivider, the conceptbox/cautionbox/examplebox/
%% takeawaybox tcolorboxes, and \compacttables used across the split decks.

\documentclass[10pt,english,aspectratio=169]{beamer}

\usepackage{ae,aecompl}
\usepackage[T1]{fontenc}
\usepackage{textcomp}
\usepackage[utf8]{inputenc}
\usepackage{babel}
\usepackage{datetime2}

\setcounter{secnumdepth}{3}
\setcounter{tocdepth}{3}

\usepackage{amsmath, amssymb, amsthm, graphicx}
\usepackage{booktabs, multirow}
\usepackage{tikz}
\usepackage{pgfplots}
\pgfplotsset{compat=1.16}
\usepackage[most]{tcolorbox}
\tcbuselibrary{skins, breakable, theorems}
\usepackage{listings}
\usepackage{xcolor}

\lstset{
  basicstyle=\ttfamily\small,
  keywordstyle=\color{blue!80!black}\bfseries,
  commentstyle=\color{green!50!black}\itshape,
  stringstyle=\color{orange!80!black},
  backgroundcolor=\color{gray!8},
  frame=single,
  rulecolor=\color{gray!40},
  breaklines=true,
  showstringspaces=false,
  numbers=none,
}

\lstdefinestyle{pythonstyle}{
    language=Python,
    basicstyle=\ttfamily\footnotesize,
    keywordstyle=\color{blue!80!black}\bfseries,
    commentstyle=\color{green!50!black}\itshape,
    stringstyle=\color{orange!80!black},
    frame=single,
    breaklines=true,
    showstringspaces=false,
    numbers=left,
    numbersep=5pt,
    numberstyle=\tiny\color{gray},
    backgroundcolor=\color{gray!8},
    rulecolor=\color{gray!40}
}

\ifx\hypersetup\undefined
  \AtBeginDocument{%
    \hypersetup{bookmarks=false,breaklinks=true,pdfborder={0 0 0},colorlinks=true,
      linkcolor=DarkText,urlcolor=RedTitle}%
  }
\else
  \hypersetup{bookmarks=false,breaklinks=true,pdfborder={0 0 0},colorlinks=true,
    linkcolor=DarkText,urlcolor=RedTitle}%
\fi

\makeatletter
\newcommand\makebeamertitle{%
  \begin{frame}[plain]
    \vspace*{1.0cm}
    \centering
    {\usebeamerfont{title}\usebeamercolor[fg]{title}\inserttitle\par}
    \vspace{1.0cm}
    {\usebeamerfont{author}\usebeamercolor[fg]{author}\insertauthor\par}
    \vspace{0.2cm}
    {\usebeamerfont{date}\usebeamercolor[fg]{date}\insertdate\par}
  \end{frame}}%
\AtBeginDocument{%
  \let\origtableofcontents=\tableofcontents
  \def\tableofcontents{\@ifnextchar[{\origtableofcontents}{\gobbletableofcontents}}
  \def\gobbletableofcontents#1{\origtableofcontents}%
}

\usetheme{metropolis}
\usefonttheme{professionalfonts}

\definecolor{LightGray}{RGB}{230,230,230}
\definecolor{RedTitle}{RGB}{0,0,200}
\definecolor{VeryLightGray}{RGB}{245,245,245}
\definecolor{DarkText}{RGB}{0,0,0}

\setbeamercolor{background canvas}{bg=white}
\setbeamercolor{title}{fg=RedTitle}
\setbeamercolor{author}{fg=DarkText}
\setbeamercolor{date}{fg=DarkText}
\setbeamercolor{frametitle}{bg=VeryLightGray,fg=DarkText}
\setbeamercolor{normal text}{fg=DarkText}
\setbeamerfont{title}{size=\large}

\setbeamersize{text margin left=5mm, text margin right=5mm}
\addtobeamertemplate{frametitle}{}{\vspace{-0.5em}}
\newcommand{\reducevspace}{\vspace{-0.5cm}}

% Shared section divider and box styles for split topic decks.
\newcommand{\sectiondivider}[2]{%
\begin{frame}[plain,noframenumbering]
\begin{center}
\vfill
{\Large\textbf{#1}\par}
\vspace{0.35cm}
{\normalsize #2\par}
\vfill
\end{center}
\end{frame}
}

\newtcolorbox{conceptbox}[1][]{
  colback=blue!3,
  colframe=RedTitle!55,
  boxrule=0.35mm,
  arc=1mm,
  left=2mm,
  right=2mm,
  top=1mm,
  bottom=1mm,
  #1
}
\newtcolorbox{cautionbox}[1][]{
  colback=red!3,
  colframe=red!50!black,
  boxrule=0.35mm,
  arc=1mm,
  left=2mm,
  right=2mm,
  top=1mm,
  bottom=1mm,
  #1
}
\newtcolorbox{examplebox}[1][]{
  colback=green!3,
  colframe=green!45!black,
  boxrule=0.35mm,
  arc=1mm,
  left=2mm,
  right=2mm,
  top=1mm,
  bottom=1mm,
  #1
}
\newtcolorbox{takeawaybox}[1][]{
  colback=VeryLightGray,
  colframe=RedTitle!55,
  boxrule=0.35mm,
  arc=1mm,
  left=2mm,
  right=2mm,
  top=1mm,
  bottom=1mm,
  #1
}
\newcommand{\compacttables}{\renewcommand{\arraystretch}{1.15}}

\setbeamertemplate{navigation symbols}{}
\setbeamertemplate{footline}{%
  \hfill%
  \usebeamercolor[fg]{page number in head/foot}%
  \usebeamerfont{page number in head/foot}%
  \insertframenumber\,/\,\inserttotalframenumber\kern1em\vskip2pt%
}

\makeatother

\author{Zhigang Feng}
% "date with time": compile timestamp so each build is identifiable.
\date{\today\ \DTMcurrenttime}


\graphicspath{{./pic/}{../../MiniCourse_8hr/Slides/Module3_LLM_Text/pic/}{../}}

\title[AI \& ML for Quant Macro]{Machine Learning and Agentic AI for Macroeconomic Research\\
Lecture 7 --- Topic 7.1: Why LLMs, Tokenization, and Word Embeddings}

\begin{document}

\makebeamertitle

%=============================================================================
\begin{frame}{Topic 7.1 Agenda}
\begin{enumerate}
  \item \textbf{Why LLMs? Motivation and History}
  \medskip
  \item \textbf{Tokenization} --- BPE, WordPiece, SentencePiece
  \medskip
  \item \textbf{Word Embeddings} --- TF-IDF $\to$ Word2Vec $\to$ contextual
\end{enumerate}
\end{frame}

%=============================================================================
\section{Why LLMs? Motivation and History}

\begin{frame}{Why Economists Care About LLMs}
\begin{itemize}
\item \textbf{Text as data} --- modern economic research increasingly leverages textual sources:
\begin{itemize}
    \item Corporate filings (10-K, 10-Q, MD\&A)
    \item Central bank speeches, FOMC minutes, ECB statements
    \item Earnings calls, analyst reports, regulatory filings
    \item News, social media, congressional records
\end{itemize}

\medskip
\item \textbf{Applications:}
\begin{itemize}
    \item Sentiment analysis for financial forecasting
    \item Policy stance detection and uncertainty indices
    \item New text-derived macroeconomic indicators
    \item Summarization of long, technical documents
\end{itemize}

\medskip
\item \textbf{Why LLMs in particular:}
\begin{itemize}
    \item Capture nuanced language and domain-specific jargon
    \item Handle long-range context (sentences, paragraphs, whole documents)
    \item Provide a unified architecture for many tasks
\end{itemize}
\end{itemize}
\end{frame}

\begin{frame}{Three Roles of Text in Economic Research}
\begin{center}
\renewcommand{\arraystretch}{1.35}
\begin{tabular}{l p{4cm} p{6cm}}
\toprule
\textbf{Role} & \textbf{What text is} & \textbf{Example} \\
\midrule
Measurement tool & A way to construct previously unobserved variables & Baker--Bloom--Davis (2016) Economic Policy Uncertainty index from newspapers \\
Behavioral outcome & A response variable revealing decisions or institutional change & Hansen, McMahon \& Prat (2018): how FOMC transparency reform changed committee speech \\
Treatment variable & A quasi-experimental input affecting behavior & Central bank communication shocks and inflation expectations \\
\bottomrule
\end{tabular}
\end{center}

\medskip
\textbf{Take-away:} the same text can play very different econometric roles. Choose your representation method to match the role---transparency matters for measurement, while subtle context matters for behavior modeling.
\end{frame}

\begin{frame}{A Brief History: RNN $\rightarrow$ Transformer $\rightarrow$ GPT}
\begin{itemize}
\item \textbf{2003:} Bengio et al. neural language model --- continuous word vectors
\medskip
\item \textbf{2013:} Mikolov et al. \emph{Word2Vec} --- dense word embeddings at scale
\medskip
\item \textbf{2014:} Pennington et al. \emph{GloVe} --- global co-occurrence matrices
\medskip
\item \textbf{2017:} Vaswani et al. \emph{``Attention Is All You Need''} --- the Transformer replaces RNNs/LSTMs
\medskip
\item \textbf{2018:} Devlin et al. \emph{BERT} (encoder); Radford et al. \emph{GPT} (decoder)
\medskip
\item \textbf{2020--2022:} GPT-3, instruction tuning, RLHF $\rightarrow$ ChatGPT moment
\medskip
\item \textbf{2023--2026:} GPT-4/4o, Claude (3, 4), Gemini, LLaMA-3, DeepSeek; long-context (100k--1M tokens), open-weights ecosystem, agentic systems
\end{itemize}

\medskip
\textbf{Driver of the shift:} attention scales better than recurrence, both in compute (parallel) and in quality (long-range dependencies).
\end{frame}

\begin{frame}{The Autoregressive Principle}
\begin{itemize}
\item An LLM is fundamentally a \textbf{next-token predictor}:
\[
   p\bigl(t^{(n+1)} \mid t^{(1)}, t^{(2)}, \ldots, t^{(n)}\bigr).
\]

\medskip
\item \textbf{Generation:} sample $t^{(n+1)}$ from this distribution, append to the context, and repeat.

\medskip
\item \textbf{What this single equation enables} (we preview the intuition on the next two slides; a full revisit comes after we understand pretraining and fine-tuning):
\begin{itemize}
    \item \textbf{Summarizing 10-K filings or FOMC minutes} --- prompt with document + instruction; the model predicts summary tokens one by one
    \item \textbf{Generating narrative forecasts} conditional on macro state --- condition on data in the prompt; next tokens extend the narrative
    \item \textbf{Producing scenario analyses in plain language} --- the scenario is the context; prediction fills in the implications
    \item \textbf{Writing and debugging code} for empirical work --- same objective over a code corpus yields code completions
\end{itemize}

\medskip
\item \textbf{Key insight:} ``everything is next-token prediction'' is mathematically tiny but pedagogically and practically enormous.
\end{itemize}
\end{frame}

\begin{frame}{Intuition: Summarizing a 10-K Filing}
\begin{itemize}
\item \textbf{Task:} condense 50 pages of an SEC 10-K risk section into 3 readable bullets.

\medskip
\item \textbf{How the autoregressive equation achieves this:}
\begin{enumerate}
    \item Concatenate the instruction and document into the prompt:
    \begin{center}
    \small\texttt{``Summarize the following 10-K risk section in 3 bullets:\textbackslash n\textbackslash n[text]\textbackslash n\textbackslash nSummary:''}
    \end{center}
    \item The model predicts token $t^{(n+1)}$ from the full prompt; this becomes the first word of the summary
    \item Each subsequent token is predicted given everything before it (including the already-generated summary words)
    \item Generation stops at a period or end-of-sequence marker
\end{enumerate}

\medskip
\item \textbf{Why does it work?} During pretraining the model read millions of (long document, short summary) pairs. The next-token distribution has implicitly learned: \emph{``after an instruction like this, produce a concise factual summary.''}

\medskip
\item \textbf{Economic punchline:} a 50-page MD\&A section is compressed into 3 bullets in under one second --- and the same pipeline scales across thousands of filings.
\end{itemize}
\end{frame}

\begin{frame}{Intuition: Extracting Stance from FOMC Minutes}
\begin{itemize}
\item \textbf{Task:} classify each sentence in FOMC minutes as \emph{hawkish}, \emph{dovish}, or \emph{neutral}.

\medskip
\item \textbf{How the autoregressive equation achieves this:}
\begin{enumerate}
    \item Prompt: \texttt{``Sentence: [sentence]\textbackslash nStance (hawkish/dovish/neutral):''}
    \item The model predicts the next token; because the vocabulary contains these exact words, the top-probability token \emph{is} the label
    \item With fine-tuning on labeled FOMC data, the next-token distribution concentrates on monetary-policy-relevant labels
\end{enumerate}

\medskip
\item \textbf{What fine-tuning changes:} the base model already understands ``hawkish'' and ``dovish''. Fine-tuning shifts the distribution so that, after a central-bank sentence, stance-label tokens are consistently most probable.

\medskip
\item \textbf{Economic application:} automated stance extraction across decades of FOMC minutes --- enabling the kind of large-scale communication analysis in Hansen, McMahon \& Prat (2018) without manual coding.
\end{itemize}
\end{frame}

\begin{frame}{An Identification-First View of Text}
\begin{itemize}
\item \textbf{Engineering perspective first:} to a machine, a document is just a character string. But every string was \emph{produced} --- by a person, in a situation, for an audience. Before any econometrics, ask: what generated this string?

\medskip
\item \textbf{The data-generating process (DGP):}
\[
   x_i \;=\; g\bigl(s_i,\; a_i,\; c_i\bigr) \;+\; \varepsilon_i
\]
\begin{itemize}
    \item $s_i$: the latent \textbf{economic state} we want to measure (policy stance, sentiment, firm risk)
    \item $a_i$: the speaker's \textbf{strategic framing} --- audience management, spin, emphasis, tone chosen for effect
    \item $c_i$: \textbf{institutional context} --- venue, legal disclosure requirements, format conventions
\end{itemize}

\medskip
\item \textbf{Why this matters for a downstream regression} $y_i = \beta\,m(x_i) + \gamma z_i + u_i$:
\begin{itemize}
    \item If framing $a_i$ is correlated with the error $u_i$, the text feature $m(x_i)$ is \emph{endogenous} --- the same omitted-variable problem as always, but hidden inside the language
    \item Measurement error is non-classical: it correlates with the true state $s_i$, biasing $\hat\beta$
\end{itemize}

\medskip
\item \textbf{Two research uses:} (a)~\textbf{Measuring} $s_i$ --- choose representations that filter out $a_i$; (b)~\textbf{Studying strategic communication} --- $a_i$ is your outcome variable (e.g.\ Hansen, McMahon \& Prat 2018).
\end{itemize}
\end{frame}

%=============================================================================
\section{Tokenization}

\begin{frame}{Tokenization: The First Step}
\begin{itemize}
\item \textbf{Goal:} convert raw character sequences into discrete \emph{tokens} that a neural network can consume.

\medskip
\item \textbf{Why not just split by spaces?}
\begin{itemize}
    \item Many economically relevant terms combine words, numbers, and symbols (e.g.\ \texttt{S\&P500}, \texttt{CO2}, \texttt{Q4\_2025}, \texttt{stagflation})
    \item Punctuation and formatting carry meaning (`\$1.2B' vs `1.2 billion')
    \item Different languages have different word boundaries (or none at all)
\end{itemize}

\medskip
\item \textbf{What is a token?} A flexible unit:
\begin{itemize}
    \item A whole word (`inflation')
    \item A subword (`infl', `ation')
    \item A single character or byte
    \item A special marker (\texttt{<EOS>}, \texttt{<PAD>}, \texttt{<USER>})
\end{itemize}
\end{itemize}
\end{frame}

\begin{frame}{Vocabulary and the Tokenization Map}
\begin{itemize}
\item Let $A$ be the alphabet (Unicode characters). A tokenizer is a function:
\[
   \mathcal{K}\colon A^{*} \;\longrightarrow\; \{1, 2, \ldots, M\}^{*}
\]
mapping any text string into a sequence of \textbf{numbers} --- each number is a positive integer that serves as an index into a lookup table of size $M$. The fact that the index is an integer is incidental; what matters is that it \emph{uniquely identifies} each token, the same way a ZIP code uniquely identifies a location.

\medskip
\item \textbf{Why integers?} Neural networks need numeric inputs; each token ID will index into an embedding lookup table.

\medskip
\item \textbf{Practical constraint --- context length:}
\begin{itemize}
    \item Models have a maximum input size $n$ (typically $2{,}048$ to $128{,}000$ tokens; some $1$M)
    \item Long policy documents, FOMC transcripts, or 10-K filings often exceed $n$
    \item Forces \emph{chunking}, \emph{summarization}, or \emph{retrieval} (a hint at Lec08)
\end{itemize}

\medskip
\item Vocabulary sizes are typically $30{,}000$--$100{,}000$ subword tokens. \textbf{To put this in perspective:} Merriam-Webster Unabridged lists $\approx470{,}000$ entries; a college-educated native speaker actively uses roughly 20{,}000--40{,}000 words. A subword vocabulary of $M=50{,}000$ therefore spans roughly the active vocabulary of a well-read person --- but because it tokenizes \emph{sub}words (stems, suffixes, common fragments), it effectively covers far more word forms and technical terms than a same-size word-level vocabulary would.
\end{itemize}
\end{frame}

\begin{frame}{Illustrative Example: ``GDP rose by 2\%''}
\begin{itemize}
\item A simple tokenization might produce
\[
   \text{(GDP},\; \text{rose},\; \text{by},\; \text{2},\; \%\text{)}
\]
mapping to integer IDs, e.g.\ $(541,\; 823,\; 71,\; 312,\; 49)$.

\medskip
\item \textbf{Why integer mapping matters:}
\begin{itemize}
    \item Each ID becomes the index of a learned embedding vector
    \item Numeric tokens (`2', `3', `5') get \emph{their own} embeddings
\end{itemize}

\medskip
\item \textbf{Injectivity matters for forecasting:} a well-designed tokenizer maps
\begin{itemize}
    \item ``inflation rate is 3.5\%''
    \item ``inflation rate is 3.6\%''
\end{itemize}
to \emph{different} token sequences. Otherwise small but economically meaningful differences would collapse into the same input.

\medskip
\item \textbf{Take-away:} tokenization is the silent foundation. Bad tokenization $\Rightarrow$ everything downstream is fighting noise.
\end{itemize}
\end{frame}

\begin{frame}{Subword Tokenization (BPE / WordPiece)}
\begin{itemize}
\item Modern LLMs use \textbf{subword} tokenizers, learned from a large corpus by greedy frequency merging:
\begin{itemize}
    \item Byte-Pair Encoding (BPE): GPT, LLaMA
    \item WordPiece: BERT
    \item SentencePiece: T5, multilingual models
\end{itemize}

\medskip
\item \textbf{Why subwords?}
\begin{itemize}
    \item Handle out-of-vocabulary words (\texttt{QuantitativeEasing} $\rightarrow$ \texttt{Quant}+\texttt{itative}+\texttt{Easing})
    \item Compress common words to one token; split rare or compound words
    \item Avoid the explosion of a full-word vocabulary
    \item Work across languages with the same vocabulary
\end{itemize}

\medskip
\item \textbf{Cost:} a previously coined term like \texttt{nowcasting} may split into 3 tokens; a numeric figure like \texttt{1.2345} may split into several digit-tokens, making numerical reasoning fragile.
\end{itemize}
\end{frame}

\begin{frame}{Tokenization Is a Preprocessing Choice}
\begin{itemize}
\item Tokenization sits inside a larger \textbf{preprocessing pipeline}:
\[
   x^{\text{raw}} \;\xrightarrow{\;\text{clean}\;}\; x^{\text{clean}}
   \;\xrightarrow{\;\text{tokenize}\;}\; \{t_{ij}\}
   \;\xrightarrow{\;\text{represent}\;}\; z_i,
\]
from raw scraped text $\to$ normalized text $\to$ integer token IDs $\to$ downstream numeric representation (TF-IDF, embedding, $\ldots$).

\medskip
\item \textbf{Specification curve problem} (Denny \& Spirling, 2018): just as economists report results under alternative lag orders, control sets, and sample restrictions, NLP results should be reported under alternative preprocessing pipelines.
\begin{itemize}
    \item LDA topics under different pipelines vary substantially --- often by as much as the substantive variation of interest
    \item Stemming and punctuation handling have the largest individual effects
    \item \textbf{Lesson:} the pipeline is a modeling decision; report the full range as you would a VAR robustness table
\end{itemize}

\medskip
\item \textbf{Lesson for empirical economists:} report the \emph{universe} of preprocessing choices. The ``convenient'' specification is not a result --- it is a choice. This applies equally to tokenizer choice and vocabulary size.
\end{itemize}
\end{frame}

\begin{frame}{The Seven Preprocessing Choices}
\small
\begin{itemize}
\item \textbf{Lowercasing} --- map \texttt{Fed} and \texttt{fed} to the same token. Compresses vocabulary, but erases proper-noun signal.
\item \textbf{Stemming / lemmatization} --- cut suffixes so \emph{banking/banks/banked} collapse to \emph{bank}. Can erase tense, plurality, policy-language nuance.
\item \textbf{Stop-word removal} --- drop high-frequency function words. Reduces noise; but negation (\emph{not}) and modality (\emph{may, will}) can matter for central-bank text.
\item \textbf{Punctuation handling} --- keep, strip, or tokenize separately. Percent signs, decimals, sentence boundaries all matter.
\item \textbf{$n$-gram inclusion} --- treat \emph{``yield curve''} as one token rather than two. Captures multi-word economic concepts at the cost of a larger vocabulary.
\item \textbf{Frequency cutoffs} --- drop terms in $<k$ or $>K$ documents. Removes typos and boilerplate; can also remove rare-but-informative events.
\item \textbf{Language filtering} --- restrict to English or apply per-language pipelines. Crucial for multilingual corpora like ECB speeches.
\end{itemize}

\smallskip
Each choice is an independent on/off switch $\Rightarrow$ $2^{7} = 128$ pipelines $\Rightarrow$ report a \emph{specification curve}, not a single run.
\end{frame}

%=============================================================================
\section{Word Embeddings}

\begin{frame}{Why Embed Tokens into Vectors?}
\begin{itemize}
\item Neural networks operate on real-valued vectors, not discrete symbols.

\medskip
\item An \textbf{embedding} is a learned function
\[
   \mathcal{E}\colon \{1, \ldots, M\} \;\longrightarrow\; \mathbb{R}^{d},
\]
mapping each token ID to a dense vector of dimension $d$ (typically $d=128$ to $4096$).

\medskip
\item \textbf{Think of it as a basis problem from linear algebra.} We want to find a coordinate system --- a set of $d$ orthogonal ``semantic axes'' (e.g.\ monetary-policy stance, positive valence, corporate context) --- such that every token maps to a point in $\mathbb{R}^{d}$. Once tokens are coordinates, we can do arithmetic: measure distances, compute cosine similarities, add and subtract meaning. We will later connect this to Arrow--Debreu securities in financial markets. The challenge: we do not know the right axes in advance. They must be \emph{learned} from data.

\medskip
\item \textbf{What we want from embeddings:}
\begin{itemize}
    \item Semantically similar tokens should have nearby vectors (close coordinates in the learned basis)
    \item Algebraic operations (sums, differences) should reflect semantic relations
    \item The mapping should be \emph{learned} from data, not hand-designed
\end{itemize}

\medskip
\item \textbf{Bonus:} we will also need a \emph{positional} signal so the model knows token order---more on this when we meet the Transformer.
\end{itemize}
\end{frame}

\begin{frame}{Three Approaches to Word Representation}
\textbf{The problem from the previous slide:} find a mapping from tokens to vectors that captures semantic relationships. Three generations of methods have tried to solve this, each fixing a specific failure of its predecessor. Understanding this progression explains \emph{why} the Transformer was needed.
\medskip
\begin{center}
\renewcommand{\arraystretch}{1.3}
\begin{tabular}{l l l l}
\toprule
\textbf{Method} & \textbf{Type} & \textbf{Idea} & \textbf{Year} \\
\midrule
TF-IDF & Sparse, frequency & Re-weight word counts & 1970s \\
Word2Vec & Dense, predictive & Predict neighbors in window & 2013 \\
GloVe & Dense, statistical & Factor co-occurrence matrix & 2014 \\
Transformer embeddings & Dense, contextual & Vary by surrounding tokens & 2017+ \\
\bottomrule
\end{tabular}
\end{center}

\medskip
\begin{itemize}
\item TF-IDF: a bag-of-words feature. Transparent and easy to interpret.
\item Word2Vec / GloVe: produce a single vector per word \emph{type} (\emph{static} embeddings).
\item Transformer embeddings: produce a different vector for each \emph{occurrence} of a word, depending on its surrounding context (\emph{contextual} embeddings).
\end{itemize}

\medskip
We'll meet each in turn, then see why static embeddings break down and motivate the Transformer.

\smallskip
\textbf{The progression:} TF-IDF (no learning, no context) $\to$ Word2Vec/GloVe (learned, static per word type) $\to$ Transformer embeddings (learned, context-dependent per occurrence). Each step addresses a specific failure of the one before, and connects directly to the Limits of Static Embeddings slide later.
\end{frame}

\begin{frame}{Term Frequency (TF)}
\begin{itemize}
\item \textbf{Definition:} share of document $d$ consisting of word $w$:
\[
   \text{TF}(w, d) \;=\; \frac{\text{count}(w, d)}{\text{total words in } d}.
\]

\medskip
\item \textbf{Toy example.} FOMC-like sentence $d$: \emph{``the Fed raised rates because inflation persisted''} (7 words). Every word appears once, so TF is $1/7 \approx 0.14$ for each:
\begin{center}
\begin{tabular}{l c c}
\toprule
word & count in $d$ & TF$(w,d)$ \\
\midrule
\emph{the} & 1 & $0.14$ \\
\emph{Fed} & 1 & $0.14$ \\
\emph{raised} & 1 & $0.14$ \\
\emph{rates} & 1 & $0.14$ \\
\emph{inflation} & 1 & $0.14$ \\
\bottomrule
\end{tabular}
\end{center}

\medskip
\item \textbf{Limitation:} TF rewards \emph{frequent} words. Across a corpus, function words (\emph{the, a, is}) dominate even though they carry almost no semantic signal. We need a counterweight.
\end{itemize}
\end{frame}

\begin{frame}{Inverse Document Frequency (IDF)}
\begin{itemize}
\item \textbf{Definition:} how rare word $w$ is across a collection of $N$ documents:
\[
   \text{IDF}(w) \;=\; \log\!\left(\frac{N}{\#\{\text{docs containing } w\}}\right).
\]

\medskip
\item \textbf{Toy example.} Corpus of $N = 100$ FOMC speeches:
\begin{center}
\begin{tabular}{l c c}
\toprule
word & docs containing it & IDF$(w)$ \\
\midrule
\emph{the}         & 100 & $\log(100/100) = 0$ \\
\emph{inflation}   & 60  & $\log(100/60) \approx 0.51$ \\
\emph{yield curve} & 5   & $\log(100/5) \approx 3.00$ \\
\bottomrule
\end{tabular}
\end{center}

\medskip
\item \textbf{Intuition:} rare $\Rightarrow$ informative $\Rightarrow$ up-weighted. \emph{``the''} gets $0$ automatically --- no stop-word list needed. \emph{``yield curve''} gets a large multiplier because it distinguishes a few speeches from the rest.
\end{itemize}
\end{frame}

\begin{frame}{Putting It Together: TF-IDF}
\begin{itemize}
\item \textbf{Product formula:} each (word, document) pair gets a weighted score
\[
   \text{TF-IDF}(w, d) \;=\; \text{TF}(w, d)\,\cdot\,\text{IDF}(w).
\]
Stacking across all words in vocabulary $V$ $\Rightarrow$ each document is a sparse vector in $\mathbb{R}^{|V|}$.

\medskip
\item \textbf{Toy 3-speech matrix} (illustrative scores):
\begin{center}
\begin{tabular}{l c c c}
\toprule
 & \emph{the} & \emph{inflation} & \emph{yield curve} \\
\midrule
speech A (hawkish)   & 0 & 0.08 & 0.00 \\
speech B (dovish)    & 0 & 0.05 & 0.00 \\
speech C (technical) & 0 & 0.02 & 0.21 \\
\bottomrule
\end{tabular}
\end{center}

\medskip
\item \textbf{Properties:}
\begin{itemize}
    \item Down-weights common words automatically (no stop-word list)
    \item Highlights distinctive words for each document
    \item Lives in $\mathbb{R}^{|V|}$ --- very high-dimensional, very sparse
    \item \emph{No} semantic similarity: \emph{``central bank''} and \emph{``Federal Reserve''} are orthogonal $\Rightarrow$ motivation for dense embeddings (next slide).
\end{itemize}
\end{itemize}
\end{frame}

\begin{frame}{Word2Vec: Dense Embeddings from Context}
\begin{itemize}
\item \textbf{Distributional hypothesis} (Firth, 1957): \emph{``You shall know a word by the company it keeps.''}

\medskip
\item \textbf{Idea} (Mikolov et al., 2013): for each word in a corpus, look at a window of $\pm L$ neighboring words. Learn a vector for the center word that \emph{predicts} its neighbors (\emph{Skip-gram}) or learn one that \emph{is predicted by} its neighbors (\emph{CBOW}).

\medskip
\item \textbf{Training output:} a single static vector $\mathbf{u}_w \in \mathbb{R}^{d}$ for each word $w$ in the vocabulary, with $d$ typically $100$--$300$.

\medskip
\item \textbf{Geometric magic:} after training, simple vector arithmetic captures analogies, e.g.
\[
   \mathbf{u}_{\text{king}} - \mathbf{u}_{\text{man}} + \mathbf{u}_{\text{woman}} \;\approx\; \mathbf{u}_{\text{queen}}.
\]
Words living on the same semantic axis differ by approximately the same vector.
\end{itemize}
\end{frame}

\begin{frame}{Skip-Gram in One Slide: The Math}
\begin{itemize}
\item Each word $w$ has \emph{two} learned vectors: an input vector $\mathbf{u}_{w}$ and an output vector $\mathbf{v}_{w}$ (often merged after training).

\medskip
\item For a center word $w_t$ and a context word $w_c$ in its window, the model predicts
\[
   P\bigl(w_c \mid w_t\bigr) \;=\; \frac{\exp\bigl(\mathbf{u}_{w_t}^{\top} \mathbf{v}_{w_c}\bigr)}{\sum_{j=1}^{M}\exp\bigl(\mathbf{u}_{w_t}^{\top} \mathbf{v}_{w_j}\bigr)}.
\]

\medskip
\item Training objective (negative log-likelihood over all (center, context) pairs):
\[
   \mathcal{L} \;=\; -\sum_{t}\sum_{c \in \text{window}(t)} \log P\bigl(w_c \mid w_t\bigr).
\]

\medskip
\item Optimized with stochastic gradient descent (in practice with \emph{negative sampling} to avoid the full vocabulary softmax).

\medskip
\item \textbf{Result:} words with similar context distributions end up with similar $\mathbf{u}$ vectors.
\end{itemize}
\end{frame}

\begin{frame}{Skip-Gram: Walkthrough on ``central banks target inflation''}
\begin{itemize}
\item \textbf{Sentence:} \emph{``central banks target inflation''}, window radius $L = 2$, center word $=$ \emph{target}.

\medskip
\item Training pairs formed from this center word:
\[
(\textit{target}, \textit{central}), \quad (\textit{target}, \textit{banks}), \quad (\textit{target}, \textit{inflation}).
\]

\medskip
\item For each pair, the score is the dot product $\mathbf{u}_{\text{target}}^{\top} \mathbf{v}_{\text{context}}$. Gradient descent \emph{increases} this score for observed pairs and \emph{decreases} it for random (negative) pairs.

\medskip
\item \textbf{Across the whole corpus:} \emph{target} appears near \emph{rates}, \emph{inflation}, \emph{policy}, \emph{Fed}. So do \emph{raise} and \emph{cut}. Consequently
\[
\mathbf{u}_{\text{target}} \;\text{ends up close to}\; \mathbf{u}_{\text{raise}},\, \mathbf{u}_{\text{cut}} \;\text{and far from}\; \mathbf{u}_{\text{apple}}.
\]

\medskip
\item \textbf{Takeaway:} analogies emerge because consistent co-occurrences place related words at similar positions in $\mathbb{R}^{d}$.
\end{itemize}
\end{frame}

\begin{frame}{GloVe: Global Co-occurrence Statistics}
\begin{itemize}
\item Pennington et al.\ (2014): instead of sliding-window prediction, factor a \emph{global} matrix.

\medskip
\item Let $X_{ij}$ count how often word $j$ appears near word $i$ across the entire corpus. GloVe fits embeddings to satisfy
\[
   \mathbf{u}_{i}^{\top}\,\mathbf{v}_{j} + b_{i} + b_{j} \;\approx\; \log X_{ij}.
\]

\medskip
\item \textbf{Toy counts on a central-bank corpus} $X_{ij}$ (illustrative):
\begin{center}
\small
\renewcommand{\arraystretch}{1.1}
\begin{tabular}{l c c c}
\toprule
 & \emph{rates} & \emph{inflation} & \emph{apple} \\
\midrule
Fed      & 820 & 410 & 3 \\
ECB      & 540 & 390 & 2 \\
hawkish  & 210 & 160 & 0 \\
\bottomrule
\end{tabular}
\end{center}
$\log X_{\text{Fed},\text{rates}} \approx 6.7$ (large target dot-product) vs.\ $\log X_{\text{Fed},\text{apple}} \approx 1.1$ (small). Fitting the equation therefore forces $\mathbf{u}_{\text{Fed}}$ close to $\mathbf{v}_{\text{rates}}, \mathbf{v}_{\text{inflation}}$ and orthogonal to $\mathbf{v}_{\text{apple}}$ --- monetary-policy structure falls out of corpus-wide statistics.

\medskip
\item \textbf{Contrast with Word2Vec:} predictive / local / streaming (W2V) vs.\ statistical / global / one-shot factorization (GloVe). Both produce \emph{static} dense vectors with similar geometry. Pick W2V for massive streaming corpora, GloVe when a full co-occurrence matrix is cheap to build.
\end{itemize}
\end{frame}

\begin{frame}{A 3D Toy of Word Embeddings}
\begin{center}
\includegraphics[width=10cm,height=5.8cm]{word_examples.png}
\par\end{center}
\small
The three axes loosely correspond to ``Animal-ness'', ``Royalty'', and ``Fruit-ness''. Words like \emph{cat}, \emph{king}, and \emph{apple} cluster along the dimension that fits them best. Real LLM embeddings live in hundreds of dimensions, but the same intuition holds.
\end{frame}

\begin{frame}{Analogy: Embeddings as Arrow--Debreu Securities}
\begin{itemize}
\item For an economist, a $d$-dimensional embedding is reminiscent of an asset's payoff vector across $d$ states of the world.

\medskip
\item \textbf{Each dimension} = a latent ``semantic state'' (royalty, animal-ness, monetary-policy stance, $\ldots$).

\medskip
\item \textbf{Each word} = a vector of exposures to those states.

\medskip
\item \textbf{Dot product} = a kind of correlation:
\begin{itemize}
    \item High dot product $\Rightarrow$ similar payoff structure $\Rightarrow$ similar meaning
    \item Orthogonal vectors $\Rightarrow$ unrelated concepts
\end{itemize}

\medskip
\item \textbf{Cosine similarity} normalizes for vector length:
\[
   \cos(\mathbf{u}, \mathbf{v}) \;=\; \frac{\mathbf{u}^{\top}\mathbf{v}}{\|\mathbf{u}\|\,\|\mathbf{v}\|}.
\]

\medskip
\item This view makes embeddings feel like a familiar object: a state-contingent payoff vector in a learned semantic state space.
\end{itemize}
\end{frame}

\begin{frame}{Limits of Static Embeddings}
\begin{columns}[T]
\column{0.48\textwidth}
\begin{tcolorbox}[colback=red!5, colframe=red!50!black, title=\textbf{What static embeddings miss}]
\begin{itemize}
\item \textbf{Polysemy:} ``rate'' = interest rate? exchange rate? growth rate? One vector for all senses.
\item \textbf{Negation:} ``GDP growth NOT slowing'' and ``GDP growth slowing'' have nearly identical mean vectors.
\item \textbf{Long-range context:} fixed window misses dependencies across paragraphs.
\item \textbf{Corpus dependence:} ``mask'' before vs.\ after 2020 are very different concepts.
\end{itemize}
\end{tcolorbox}

\column{0.48\textwidth}
\begin{tcolorbox}[colback=green!5, colframe=green!40!black, title=\textbf{What we need next}]
\begin{itemize}
\item Vectors that change with \emph{context}, not just word identity
\item A mechanism for tokens to ``look at'' other tokens, near or far
\item A way to learn arbitrary relational patterns, not just co-occurrence within a window
\end{itemize}
\medskip
$\Longrightarrow$ \textbf{The Transformer.}
\end{tcolorbox}
\end{columns}
\end{frame}

\end{document}
